% =====================================================================
% Chapitre : Phase B - Étape 5 (Segmentation réseau) et Étape 6
% (Règles de pare-feu Zero Trust)
% =====================================================================

\chapter{Segmentation réseau et politique de pare-feu (Zero Trust)}

\section{Introduction}

Après avoir démontré en phase A qu'une architecture réseau plate expose
l'ensemble des services à toute machine compromise sur le réseau (voir
chapitre précédent, chaîne d'attaque complète depuis Kali jusqu'à
l'exfiltration des identifiants Nextcloud), la phase B du projet consiste
à reconstruire l'infrastructure selon les principes Zero Trust. La première
brique de cette reconstruction est la \textbf{segmentation réseau} : au lieu
d'un réseau unique où toutes les machines peuvent se joindre librement,
le réseau est désormais découpé en six segments isolés, chacun avec un rôle
précis, reliés uniquement via un pare-feu central (OPNsense) qui applique
une politique de moindre privilège.

Ce chapitre documente deux étapes indissociables : l'\textbf{étape 5}
(conception et mise en place de la segmentation : interfaces virtuelles,
adressage, migration des machines) et l'\textbf{étape 6} (définition et
implémentation des règles de pare-feu qui donnent corps à cette segmentation).
Ces deux étapes sont en réalité fortement couplées : une segmentation réseau
sans règles de filtrage précises n'apporte aucune sécurité, un point que ce
chapitre met explicitement en évidence à travers une erreur de configuration
réelle rencontrée durant le projet.

\section{Étape 5 — Segmentation réseau}

\subsection{Conception de l'architecture segmentée}

Le réseau plat de la phase A (un unique sous-réseau \texttt{192.168.1.0/24})
est remplacé par six segments distincts, chacun représenté par une
interface dédiée du pare-feu OPNsense. Le tableau~\ref{tab:segments}
résume l'architecture cible.

\begin{table}[h]
\centering
\caption{Plan d'adressage des segments réseau}
\label{tab:segments}
\begin{tabular}{|l|l|l|l|}
\hline
\textbf{Segment} & \textbf{Interface OPNsense} & \textbf{Sous-réseau} & \textbf{Rôle} \\
\hline
MGMT & LAN (em1) & 10.10.1.0/24 & Administration du lab \\
IAM / DMZ & OPT1 (em2) & 10.10.2.0/24 & Identité (Keycloak, Authelia, Traefik) \\
APPS & OPT2 (em3) & 10.10.3.0/24 & Applications (Server-Nextcloud) \\
DB & OPT3 (em4) & 10.10.4.0/24 & Base de données (Server-DB) \\
ATTACKER & OPT4 (em5) & 10.10.5.0/24 & Poste attaquant (Kali) \\
MONITORING & OPT5 (em6) & 10.10.6.0/24 & Supervision (Wazuh) \\
\hline
\end{tabular}
\end{table}

Cette séparation matérialise un principe central du Zero Trust : la
confiance n'est plus implicite du simple fait d'être « sur le réseau
interne ». Chaque segment ne doit pouvoir communiquer qu'avec les segments
strictement nécessaires à sa fonction, et dans un seul sens quand c'est
pertinent (voir \ref{sec:lecon-lateral}).

\subsection{Limitation de l'interface graphique VirtualBox et contournement}

La machine virtuelle du pare-feu (\texttt{OPNsense-FW}) nécessite sept
interfaces réseau : une interface WAN (connexion à la machine hôte) et six
interfaces internes, une par segment. Or l'interface graphique de
VirtualBox n'affiche que quatre onglets d'adaptateurs réseau par VM, même
après avoir sélectionné un chipset \texttt{ICH9} (qui supporte pourtant
jusqu'à huit adaptateurs). Les interfaces 5 à 8 ne sont configurables que
via la ligne de commande \texttt{VBoxManage}.

\begin{lstlisting}[language=bash, caption={Ajout des adaptateurs réseau 5 à 7 via VBoxManage}]
VBoxManage modifyvm "OPNsense-FW" --nic5 intnet --intnet5 "intnet-apps"
VBoxManage modifyvm "OPNsense-FW" --nic6 intnet --intnet6 "intnet-db"
VBoxManage modifyvm "OPNsense-FW" --nic7 intnet --intnet7 "intnet-attacker"
\end{lstlisting}

Chaque réseau interne (\texttt{intnet-*}) correspond à un segment isolé :
seules les VM dont l'adaptateur réseau pointe vers le même nom de réseau
interne peuvent échanger des trames Ethernet entre elles — et uniquement
via OPNsense, qui est la seule machine connectée à tous ces réseaux
internes simultanément.

\subsection{Assignation et adressage des interfaces OPNsense}

Au premier démarrage après l'ajout des interfaces, la console OPNsense
propose de réassigner les interfaces physiques (\texttt{em0} à
\texttt{em6}) aux rôles logiques (WAN, LAN, OPT1 à OPT5). L'identification
se fait par adresse MAC, affichée pour chaque interface détectée. Le WAN
est assigné en premier (interface connectée à la machine hôte), puis LAN
et les cinq interfaces OPT dans l'ordre du plan d'adressage.

Chaque interface reçoit ensuite son adresse IP via le menu console
« Set interface IP address », reproduisant le tableau~\ref{tab:segments} :
OPNsense porte systématiquement le premier hôte de chaque sous-réseau
(ex. \texttt{10.10.3.1/24} pour APPS), qui sert de passerelle par défaut
aux machines de ce segment.

\subsection{Migration des machines existantes}

Les trois machines déjà installées en phase A (Server-DB, Server-Nextcloud,
Kali-Attacker) doivent être déplacées vers leurs nouveaux segments. Cette
migration comporte systématiquement deux volets : changer l'adaptateur
réseau virtuel de la VM (dans les paramètres VirtualBox, passer de
l'ancien réseau \texttt{intnet-lan} unique vers le nouveau réseau interne
dédié), puis reconfigurer l'adresse IP statique à l'intérieur de l'OS invité.

\begin{lstlisting}[caption={Configuration réseau statique type (\texttt{/etc/network/interfaces})}]
auto lo
iface lo inet loopback

auto eth0
iface eth0 inet static
    address 10.10.5.50/24
    gateway 10.10.5.1
    dns-nameservers 10.10.5.1
\end{lstlisting}

Cette migration, en apparence triviale, a mis en évidence plusieurs
difficultés concrètes documentées ci-dessous, car chacune illustre un
piège réaliste que rencontrerait n'importe quel administrateur effectuant
une telle migration en production.

\subsubsection{Comportement « default-deny » des nouvelles interfaces OPNsense}

Après avoir migré Server-DB vers le segment DB (\texttt{10.10.4.20/24}),
le ping vers la passerelle (\texttt{10.10.4.1}) échouait à 100\,\%, alors
que la configuration IP de la machine était correcte. La cause n'était pas
un problème de configuration côté VM, mais un comportement par défaut
d'OPNsense : contrairement à l'interface LAN historique, qui est créée
avec une règle « autoriser tout le trafic sortant » par défaut lors de
l'installation, toute nouvelle interface OPT créée ultérieurement
\textbf{ne possède aucune règle autorisée par défaut}. Le pare-feu bloque
donc tout, y compris un simple ping vers sa propre passerelle, tant
qu'aucune règle explicite n'est ajoutée sur cet onglet. L'ajout d'une
règle temporaire « Pass Any » sur l'onglet OPT3 a résolu ce point de
blocage initial (règle qui sera ensuite resserrée à l'étape 6).

\subsubsection{Le problème de démarrage circulaire de \texttt{occ} (Nextcloud)}

Une fois Server-Nextcloud migré vers \texttt{10.10.3.30/24}, il fallait
indiquer à Nextcloud la nouvelle adresse de la base de données. La méthode
recommandée par Nextcloud consiste à utiliser son outil en ligne de
commande :

\begin{lstlisting}[language=bash]
sudo -u www-data php /var/www/nextcloud/occ config:system:set dbhost --value="10.10.4.20"
\end{lstlisting}

Cette commande a échoué avec une exception Doctrine/DBAL non gérée
(\texttt{Failed to connect to the database}). La raison est un problème
d'amorçage circulaire : \texttt{occ} est lui-même une application PHP qui
doit charger l'ensemble du framework Nextcloud — y compris établir une
connexion à la base de données — avant de pouvoir exécuter \emph{n'importe
quelle} commande, y compris celle censée corriger précisément cette
connexion. Tant que l'ancienne adresse de base de données
(\texttt{192.168.1.20}, devenue injoignable) figure dans la configuration,
\texttt{occ} ne peut tout simplement pas démarrer.

La solution a consisté à contourner \texttt{occ} et éditer directement le
fichier de configuration PHP :

\begin{lstlisting}[language=bash]
sudo nano /var/www/nextcloud/config/config.php
\end{lstlisting}

\begin{lstlisting}[language=php, caption={Extrait corrigé de config.php}]
'dbhost' => '10.10.4.20',
'trusted_domains' =>
array (
    0 => 'localhost',
    1 => '10.10.3.30',
    2 => '10.10.1.50',
),
\end{lstlisting}

\subsubsection{Contrôle d'accès MariaDB lié à l'hôte d'origine}

Après correction de \texttt{config.php}, \texttt{occ status} renvoyait une
nouvelle erreur distincte :

\begin{lstlisting}
SQLSTATE[HY000] [1130] Host '10.10.3.30' is not allowed to
connect to this MariaDB server
\end{lstlisting}

Ce message illustre un point pédagogique important : la segmentation
réseau (gérée par OPNsense) et le contrôle d'accès applicatif (géré par
MariaDB) sont deux couches de sécurité \textbf{indépendantes}. Même si le
pare-feu autorise désormais le trafic entre APPS et DB, MariaDB maintient
sa propre liste d'autorisations par couple \texttt{utilisateur@hôte}. Le
compte \texttt{nextcloud} avait été créé en phase A avec un hôte autorisé
basé sur l'ancien sous-réseau (\texttt{nextcloud'@'192.168.1.\%}), qui ne
correspond plus à la nouvelle adresse de Nextcloud. C'est un exemple
concret de \emph{défense en profondeur} : même si un attaquant contourne
le pare-feu (ou si celui-ci est mal configuré), l'application elle-même
applique ses propres restrictions d'accès, et une migration réseau doit
donc être répercutée à chaque couche indépendamment.

La correction consiste à créer un compte scoped sur l'adresse exacte de
Nextcloud plutôt que de réutiliser un sous-réseau trop large, puis à
supprimer l'ancien compte :

\begin{lstlisting}[language=sql]
CREATE USER 'nextcloud'@'10.10.3.30' IDENTIFIED BY 'DB_PASSWORD_REDACTED';
GRANT ALL PRIVILEGES ON nextcloud.* TO 'nextcloud'@'10.10.3.30';
FLUSH PRIVILEGES;
DROP USER 'nextcloud'@'192.168.1.%';
\end{lstlisting}

Après cette correction, \texttt{occ status} confirme
\texttt{installed: true}.

\subsubsection{Corruption du fichier de configuration réseau sur Kali}

Lors de la migration de Kali-Attacker, une édition en place de
\texttt{/etc/network/interfaces} a fusionné accidentellement l'ancien bloc
de configuration DHCP avec le nouveau bloc statique, sans supprimer le
premier — allant jusqu'à concaténer une ligne de commentaire directement
à la suite de la ligne \texttt{iface eth0 inet dhcp}. Le service réseau
refusait alors de démarrer (\texttt{control process exited with error
code}), laissant l'interface \texttt{eth0} dans l'état \texttt{DOWN}. La
correction la plus sûre dans ce cas n'est pas une édition supplémentaire
du fichier corrompu, mais une réécriture atomique complète via un
here-document, qui garantit un fichier final cohérent :

\begin{lstlisting}[language=bash]
sudo tee /etc/network/interfaces > /dev/null << 'EOF'
source /etc/network/interfaces.d/*

auto lo
iface lo inet loopback

auto eth0
iface eth0 inet static
    address 10.10.5.50/24
    gateway 10.10.5.1
    dns-nameservers 10.10.5.1
EOF
\end{lstlisting}

\subsection{Cas particulier : le poste d'administration (client-debian)}

Le segment MGMT est administré depuis une machine virtuelle Debian
dédiée (\texttt{client-debian}, \texttt{10.10.1.50/24}), distincte des
serveurs applicatifs. Un point notable rencontré sur cette machine : son
interface \texttt{enp0s3} est marquée \texttt{unmanaged} par
NetworkManager (car définie statiquement dans
\texttt{/etc/network/interfaces}), ce qui rend les commandes
\texttt{nmcli} inopérantes pour modifier sa configuration DNS. Le fichier
\texttt{/etc/resolv.conf} restait bloqué sur l'ancienne passerelle
(\texttt{192.168.1.1}) malgré un \texttt{dns-nameservers 10.10.1.1}
correctement présent dans \texttt{/etc/network/interfaces} et plusieurs
redémarrages du service \texttt{networking}. La correction a nécessité une
édition directe de \texttt{/etc/resolv.conf} (fichier simple, non
symbolique, qui n'était régénéré par aucun service actif une fois
l'interface marquée \texttt{unmanaged}).

\section{Étape 6 — Politique de pare-feu (règles de moindre privilège)}

\subsection{Une leçon essentielle : la direction du trafic et le mouvement latéral}
\label{sec:lecon-lateral}

Une fois les quatre machines actives migrées, un test de vérification a
révélé un problème critique : des pings lancés depuis Kali-Attacker
(\texttt{10.10.5.50}) vers Server-Nextcloud (\texttt{10.10.3.30}) et
Server-DB (\texttt{10.10.4.20}) \textbf{réussissaient}, alors que la
segmentation était censée les en empêcher.

La cause provenait des règles temporaires ajoutées lors de la migration
pour résoudre le comportement default-deny (voir plus haut) : des règles
du type \texttt{Pass, Source: OPTn net, Destination: Any} avaient été
laissées en place sur chaque interface. Sur un pare-feu de type
OPNsense/pfSense, les règles sont évaluées sur l'interface par laquelle le
trafic \emph{entre} dans le pare-feu, et une destination \texttt{Any} ne
signifie pas « Internet » mais littéralement \emph{toute destination
atteignable}, y compris les segments internes voisins. Une règle pensée
pour autoriser « l'accès à Internet » autorisait donc, sans le vouloir,
le mouvement latéral complet entre tous les segments internes — annulant
de fait tout le travail de segmentation réalisé à l'étape 5.

Cette découverte est la motivation directe de la conception rigoureuse des
règles ci-dessous : chaque segment ne doit avoir que des règles
\textbf{explicites}, limitées à des destinations et des ports précis,
jamais une règle large « par commodité ».

\subsection{Principe de conception des règles}

Pour chaque segment, les règles ont été conçues en répondant à une
question simple : « de quoi cette machine a-t-elle réellement besoin pour
fonctionner, et rien de plus ? ». Le tableau~\ref{tab:regles} résume la
politique retenue pour les quatre segments actifs du lab.

\begin{table}[h]
\centering
\caption{Règles de pare-feu par segment (moindre privilège)}
\label{tab:regles}
\begin{tabular}{|l|l|l|l|}
\hline
\textbf{Segment} & \textbf{Destination autorisée} & \textbf{Port(s)} & \textbf{Justification} \\
\hline
ATTACKER & Internet (Any) & 80, 443 & Aucun accès interne \\
APPS & 10.10.4.20 (DB) & 3306 & Connexion base de données \\
APPS & Internet (Any) & 80, 443 & Mises à jour, apps Nextcloud \\
DB & \emph{aucune} & \emph{aucune} & N'initie jamais de connexion \\
MGMT & Any & 443 & Administration web (OPNsense, futurs services) \\
MGMT & Any & 22 & Administration SSH \\
MGMT & Any & 80, 53 & Internet + résolution DNS \\
\hline
\end{tabular}
\end{table}

Le cas du segment DB illustre un principe de sécurité réseau important :
une base de données n'a, par nature, jamais besoin d'initier une connexion
sortante — elle ne fait qu'écouter et répondre. L'absence totale de règle
sortante sur l'interface OPT3 n'empêche pourtant pas Server-DB de répondre
aux requêtes légitimes d'APPS, car OPNsense est un pare-feu à états
(\emph{stateful}) : le trafic de retour d'une connexion déjà autorisée en
entrée (APPS vers DB) est automatiquement accepté sans règle explicite
supplémentaire. Documenter et vérifier ce comportement a permis de
confirmer qu'une politique de filtrage strict côté DB n'introduit aucune
régression fonctionnelle.

\subsection{Implémentation et vérification, segment par segment}

Chaque segment a été vérifié immédiatement après application des règles,
en combinant un test fonctionnel positif (le service continue de
fonctionner) et un test négatif (l'accès non désiré échoue effectivement).

\begin{lstlisting}[caption={Vérification ATTACKER : aucun mouvement latéral}]
$ ping -c 3 10.10.3.30
3 packets transmitted, 0 received, 100% packet loss

$ ping -c 3 10.10.4.20
3 packets transmitted, 0 received, 100% packet loss
\end{lstlisting}

\begin{lstlisting}[caption={Vérification APPS : service fonctionnel, mouvement latéral bloqué}]
$ sudo -u www-data php /var/www/nextcloud/occ status
  - installed: true

$ ping -c 3 10.10.5.50
3 packets transmitted, 0 received, 100% packet loss
\end{lstlisting}

\begin{lstlisting}[caption={Vérification DB : aucune règle sortante, service toujours accessible}]
# Depuis Server-Nextcloud, après suppression de toute règle sortante sur OPT3 :
$ sudo -u www-data php /var/www/nextcloud/occ status
  - installed: true
\end{lstlisting}

\subsection{Un bug de configuration réaliste : la règle SSH en IPv6}

La vérification du segment MGMT a révélé un bug de configuration subtil,
suffisamment instructif pour être documenté intégralement. La règle
autorisant l'administration SSH (\texttt{MGMT -> Any SSH}) avait été créée
avec la famille d'adresses \textbf{IPv6} au lieu d'IPv4 — une valeur par
défaut du formulaire OPNsense passée inaperçue à la création. Or
l'ensemble du laboratoire fonctionne exclusivement en IPv4
(\texttt{10.10.x.x}) : cette règle ne filtrait donc jamais aucun trafic
réel, et toute tentative de connexion SSH était silencieusement rejetée
par la règle implicite \texttt{Default deny} d'OPNsense, sans que la
configuration déclarée (« une règle SSH existe ») ne corresponde au
comportement observé.

Le diagnostic a été posé grâce au journal temps réel d'OPNsense
(\emph{Firewall > Log Files > Live View}), qui a montré explicitement le
paquet SSH rejeté sur l'interface LAN par la règle \texttt{Default deny /
state violation rule} :

\begin{lstlisting}
Interface  Action  Label
LAN        block   Default deny / state violation rule
  Source: 10.10.1.50   Destination: 10.10.3.30:22
\end{lstlisting}

Ce cas illustre un principe méthodologique central pour l'administration
de pare-feu : \textbf{ne jamais faire confiance à la configuration
déclarée sans la vérifier dans les journaux de trafic réel}. Une règle
peut exister, apparaître correcte à la lecture, et pourtant ne filtrer
aucun paquet du trafic visé si un paramètre secondaire (ici la famille
d'adresses) ne correspond pas au trafic réel. La correction (changement de
\texttt{IPv6} vers \texttt{IPv4} dans le champ \emph{Address Family} de la
règle) a immédiatement résolu le problème, confirmé par un test de port :

\begin{lstlisting}
$ nc -zv -w 3 10.10.3.30 22
10.10.3.30 [10.10.3.30] 22 (ssh) open
\end{lstlisting}

\subsection{Résolution DNS pour le segment MGMT}

Un second point a nécessité une règle supplémentaire non anticipée
initialement : la résolution de noms de domaine (DNS) utilise le port 53,
distinct des ports web (80/443) déjà autorisés. Sans règle explicite pour
ce port, client-debian pouvait joindre Internet par adresse IP mais pas
par nom de domaine (\texttt{Temporary failure in name resolution}). Une
règle \texttt{MGMT -> Any, port 53 (TCP/UDP)} a été ajoutée pour corriger
ce point, rappelant qu'une politique de moindre privilège doit être
construite en identifiant \emph{tous} les protocoles réellement utilisés
par un service, DNS étant souvent oublié car invisible lors d'un usage
normal.

\section{Segments documentés pour implémentation ultérieure}

Les segments IAM/DMZ et MONITORING ne portent encore aucune machine
virtuelle à ce stade du projet (ils seront respectivement peuplés aux
étapes 7 à 9 pour Keycloak/Authelia/Traefik, et à l'étape 10 pour Wazuh).
Leur politique de filtrage a néanmoins été conçue par anticipation, afin
de garantir la cohérence de l'architecture Zero Trust dès l'installation
de ces services.

\subsection{IAM / DMZ (10.10.2.0/24)}

\begin{itemize}
    \item IAM $\rightarrow$ DB (\texttt{10.10.4.20:3306}) : Keycloak
    nécessite sa propre base de données.
    \item IAM $\rightarrow$ Internet (80/443) : mises à jour, éventuels
    appels OIDC externes.
    \item APPS $\rightarrow$ IAM (443) : Nextcloud devra rediriger
    l'authentification vers Authelia/Traefik.
    \item MGMT $\rightarrow$ IAM (443, déjà couvert par la règle MGMT
    existante).
    \item Aucune règle IAM $\rightarrow$ APPS, ATTACKER ou MONITORING.
\end{itemize}

\subsection{MONITORING (10.10.6.0/24)}

\begin{itemize}
    \item APPS, DB, IAM, ATTACKER $\rightarrow$ MONITORING, ports agent
    Wazuh (1514/1515 TCP) : les agents remontent leurs événements vers le
    serveur central.
    \item MONITORING $\rightarrow$ Internet (80/443) : mise à jour des
    règles de détection.
    \item MGMT $\rightarrow$ MONITORING (443) : accès au tableau de bord
    Wazuh.
    \item \textbf{Aucune} règle MONITORING $\rightarrow$ autres segments.
\end{itemize}

Ce dernier point est un choix de conception Zero Trust délibéré : le flux
de supervision est strictement unidirectionnel, des machines surveillées
vers le collecteur (\emph{push}), jamais l'inverse. Cela limite
drastiquement la surface d'attaque exploitable si le serveur de
supervision venait lui-même à être compromis : il ne pourrait pas servir
de rebond vers les autres segments.

\section{Synthèse}

Le tableau~\ref{tab:synthese-step6} résume l'état de la segmentation et
du filtrage à l'issue de l'étape 6.

\begin{table}[h]
\centering
\caption{État des segments à l'issue de l'étape 6}
\label{tab:synthese-step6}
\begin{tabular}{|l|l|l|}
\hline
\textbf{Segment} & \textbf{Règles de pare-feu} & \textbf{État} \\
\hline
ATTACKER & Implémentées et vérifiées & Opérationnel \\
APPS & Implémentées et vérifiées & Opérationnel \\
DB & Implémentées et vérifiées & Opérationnel \\
MGMT & Implémentées et vérifiées & Opérationnel \\
IAM / DMZ & Conçues, non implémentées & En attente (étapes 7-9) \\
MONITORING & Conçues, non implémentées & En attente (étape 10) \\
\hline
\end{tabular}
\end{table}

Cette phase a permis de transformer le réseau plat et vulnérable de la
phase A en une architecture segmentée où chaque composant ne dispose que
des accès strictement nécessaires à son fonctionnement. Elle a également
mis en évidence, à travers des incidents réels rencontrés durant la mise
en œuvre (règles trop permissives, bug de famille d'adresses IPv6, port
DNS oublié), que la segmentation réseau est un exercice itératif qui
nécessite une vérification systématique par les journaux de trafic, et
non une simple confiance dans la configuration déclarée. L'étape 12 du
projet consistera à rejouer la chaîne d'attaque de la phase A sur cette
nouvelle architecture, afin de mesurer concrètement l'efficacité de cette
segmentation face aux mêmes techniques d'attaque.
